%% toc: P2 19 19.8
\section{Equality boundaries and the connection to superposition}
\label{sec:p02-w02b-p02-026-equality}

Theorem~\ref{thm:fol-resolution} and its proved form
Theorem~\ref{thm:p02-w02-p02-028-resolution} concern equality-free clause
sets only, and the paragraph closing \cref{sec:fol-cardinality-normalforms-04}
records that bare resolution is not complete for logical equality.  This
section supplies one equality-complete connection in full: the clausal
equality axioms of the input's symbols are added to the input, and the
Resolution and Factoring rules of \cref{sec:fol-cardinality-normalforms-04}
are then used unchanged.  It also determines exactly what bare resolution
refutes when equality literals are present, and proves the soundness of the
ordered-superposition rules that replace the axioms in practice.

\emph{Purpose.}  To extend the soundness and refutational completeness of
first-order resolution from equality-free clause sets to clause sets whose
literals may contain \(=\), with \(=\) read as identity in every structure.
\emph{Inputs.}  A finite set \(S\) of first-order clauses of a signature
\(\Lang\) in the sense of \cref{sec:fol-cardinality-normalforms-04}; its
literals may be equations \(r=t\) or their negations.
\emph{Conclusions.}  Theorem~\ref{thm:p02-w02b-p02-026-eq-resolution}
(soundness and refutational completeness of resolution with the clausal
equality axioms), Theorem~\ref{thm:p02-w02b-p02-026-bare} (what resolution
without the axioms refutes), Corollary~\ref{thm:p02-w02b-p02-026-failures}
(two failures of incomplete treatments) and
Lemma~\ref{thm:p02-w02b-p02-026-superposition-sound} (soundness of the
superposition rules).  All proofs are in the annex.

\subsection{Identity semantics and the relational reading}
\label{sec:p02-w02b-p02-026-readings}

\begin{hbconvention}{def:p02-w02b-p02-026-reading}{Identity semantics, relational reading and translation}
\begin{enumerate}
  \item \emph{Identity semantics.}  As everywhere in this part, an
        \(\Lang\)-structure interprets \(=\) as the identity relation of its
        nonempty domain, and a structure satisfies a clause \(C\) iff it
        satisfies its universal closure \(\operatorname{ucl}_{\Lang}(C)\)
        (\cref{sec:fol-cardinality-normalforms-04}).  A clause set is
        \emph{satisfiable} when some \(\Lang\)-structure satisfies all its
        members, and \emph{unsatisfiable} otherwise.
  \item \emph{Ground terms.}  As in
        Convention~\ref{def:p02-w02-p02-028-standing}, \(\Lang\) contains a
        constant; when it has none, one fresh constant is added first.
  \item \emph{The fresh symbol.}  \(\mathsf{eq}\) is a binary relation
        symbol not in \(\Lang\), written infix.  Under the \emph{relational
        reading} an \((\Lang\cup\{\mathsf{eq}\})\)-structure may interpret
        \(\mathsf{eq}\) as any binary relation on its domain.
  \item \emph{Translation.}  For \(\Lang\)-atoms put
        \[
         (r=t)^{\mathsf{eq}}=(r\mathrel{\mathsf{eq}}t),
         \qquad
         R(t_1,\ldots,t_n)^{\mathsf{eq}}=R(t_1,\ldots,t_n);
        \]
        for a literal
        \((\neg A)^{\mathsf{eq}}=\neg(A^{\mathsf{eq}})\); for a clause
        \(C^{\mathsf{eq}}=\{L^{\mathsf{eq}}:L\in C\}\); for a clause set
        \(T^{\mathsf{eq}}=\{C^{\mathsf{eq}}:C\in T\}\).  Thus
        \(C^{\mathsf{eq}}\) is equality-free in the sense of
        \cref{sec:fol-cardinality-normalforms-04}, and translation leaves
        every term unchanged.
  \item \emph{Identity expansion.}  For an \(\Lang\)-structure
        \(\mathcal N\), \(\mathcal N^{=}\) is its expansion to
        \(\Lang\cup\{\mathsf{eq}\}\) that interprets \(\mathsf{eq}\) as
        identity.
\end{enumerate}
\end{hbconvention}

The resolution calculus of \cref{sec:fol-cardinality-normalforms-04} never
consults the meaning of \(=\): its unification procedure treats \(=\) as one
more binary atomic symbol, pairing the two ordered arguments.  Resolution is
therefore sound for clauses with equality (the soundness theorem of that
section holds for every structure), but it sees \(=\) only through the
relational reading.  The translation lemma makes this exact.

\begin{hblemma}{thm:p02-w02b-p02-026-translation}{proved=annex:p02-w02b-p02-026-translation}{Translation lemma}
Let \(L\) be an \(\Lang\)-literal, \(C\) an \(\Lang\)-clause, \(A,B\)
\(\Lang\)-atoms, \(\theta\) a term substitution and \(T\) a set of
\(\Lang\)-clauses.
\begin{enumerate}
  \item \((L\theta)^{\mathsf{eq}}=L^{\mathsf{eq}}\theta\),
        \((C\theta)^{\mathsf{eq}}=C^{\mathsf{eq}}\theta\) and
        \(\operatorname{Var}(C^{\mathsf{eq}})=\operatorname{Var}(C)\); the
        translation is injective on literals.
  \item \(\theta\) is a unifier of \(A\) and \(B\) iff it is a unifier of
        \(A^{\mathsf{eq}}\) and \(B^{\mathsf{eq}}\), and a most general
        unifier of the one pair iff of the other.
  \item A finite sequence \(C_1,\ldots,C_n\) of \(\Lang\)-clauses is a
        first-order resolution derivation from \(T\) iff
        \(C_1^{\mathsf{eq}},\ldots,C_n^{\mathsf{eq}}\) is one from
        \(T^{\mathsf{eq}}\).
  \item For every \(\Lang\)-structure \(\mathcal N\) and assignment \(s\),
        \(\mathcal N,s\entails L\) iff
        \(\mathcal N^{=},s\entails L^{\mathsf{eq}}\); hence \(\mathcal N\)
        satisfies \(C\) iff \(\mathcal N^{=}\) satisfies \(C^{\mathsf{eq}}\).
\end{enumerate}
\end{hblemma}

\subsection{The clausal equality axioms}
\label{sec:p02-w02b-p02-026-axioms}

\begin{hbdefinition}{def:p02-w02b-p02-026-eq-clauses}{Clausal equality axioms}
Let \(\Sigma\) be a set of constant, function and relation symbols of
\(\Lang\).  In each clause below the displayed variables are pairwise
distinct.  \(\mathrm{Eq}(\Sigma)\) is the set of the following
\(\Lang\)-clauses:
\begin{enumerate}
  \item[(E1)] reflexivity \(\{x=x\}\);
  \item[(E2)] symmetry \(\{\neg(x=y),\ y=x\}\);
  \item[(E3)] transitivity \(\{\neg(x=y),\ \neg(y=z),\ x=z\}\);
  \item[(E4)] for each \(n\)-ary function symbol \(f\in\Sigma\) with
        \(n\geq1\), function congruence
        \[
         \{\neg(x_1=y_1),\ldots,\neg(x_n=y_n),\
           f(x_1,\ldots,x_n)=f(y_1,\ldots,y_n)\};
        \]
  \item[(E5)] for each \(n\)-ary relation symbol \(R\in\Sigma\) with
        \(n\geq1\), relation congruence
        \[
         \{\neg(x_1=y_1),\ldots,\neg(x_n=y_n),\
           \neg R(x_1,\ldots,x_n),\ R(y_1,\ldots,y_n)\}.
        \]
\end{enumerate}
For a clause set \(S\), \(\Sigma_S\) is the set of constant, function and
relation symbols occurring in \(S\), and
\(\mathrm{Eq}_S\defeq\mathrm{Eq}(\Sigma_S)\).  When \(S\) is finite, so is
\(\mathrm{Eq}_S\).
\end{hbdefinition}

These are the clause forms of the congruence axioms that Nieuwenhuis and
Rubio list for treating equality by resolution, written there as
implications: reflexivity, symmetry, transitivity, one monotonicity axiom of
the first kind for each non-constant function symbol, and one of the second
kind for each predicate symbol \cite[\S1.1, p.~3]{NieuwenhuisRubio2001Typescript}.
They state that a clause set is satisfiable in first-order logic with
equality iff its union with these axioms is satisfiable without equality
\cite[\S1.1, p.~4]{NieuwenhuisRubio2001Typescript};
Lemma~\ref{thm:p02-w02b-p02-026-quotient} proves that statement for the
clause sets of this section.  The sentence form
\(\mathrm{Eq}_{\mathcal L}\) of Definition~\ref{def:p01-w02-p01-025-eq-axioms}
has the same function and relation congruence content and replaces symmetry
and transitivity by the congruence axiom for \(=\);
Theorem~\ref{thm:p01-w02-p01-025-axiom-metatheorem} is the natural-deduction
counterpart of the theorem below.  A relation symbol of arity \(0\) needs no
clause, and constants need none either: the function congruence clause of a
constant \(c\) would be \(\{c=c\}\), an instance of (E1).

\begin{hblemma}{thm:p02-w02b-p02-026-quotient}{proved=annex:p02-w02b-p02-026-quotient}{The equality axioms capture identity}
Let \(T\) be a set of \(\Lang\)-clauses and \(\Sigma\) a set of symbols of
\(\Lang\) that contains every constant, function and relation symbol
occurring in \(T\).  Then \(T\) is satisfiable iff
\(T^{\mathsf{eq}}\cup\mathrm{Eq}(\Sigma)^{\mathsf{eq}}\) is satisfied by some
\((\Lang\cup\{\mathsf{eq}\})\)-structure.  The proof of the right-to-left
direction uses no choice principle.
\end{hblemma}

\subsection{Resolution with the equality axioms}
\label{sec:p02-w02b-p02-026-method}

\begin{hbmethod}{def:p02-w02b-p02-026-method}{Resolution with the clausal equality axioms}
\emph{Input:} a finite set \(S\) of \(\Lang\)-clauses, possibly containing
equations.
\begin{enumerate}
  \item Compute \(\Sigma_S\) from the clauses of \(S\) as given, after any
        Skolemization and clause-form conversion, so that every symbol
        introduced by those steps is included.
  \item Form the finite set \(S\cup\mathrm{Eq}_S\).
  \item Build first-order resolution derivations from \(S\cup\mathrm{Eq}_S\)
        with the Resolution and Factoring rules of
        \cref{sec:fol-cardinality-normalforms-04}, with all their side
        conditions: both parents of a Resolution step are replaced by
        renamed-apart copies; the resolved literals are a positive literal
        \(A\) of the first parent and a negative literal \(\neg B\) of the
        second; Factoring merges two distinct literals of the same polarity;
        and in each case the substitution applied is a most general unifier of
        the two atoms, computed by the unification procedure of that section
        (Lemma~\ref{thm:p02-w02-p02-028-unification}).  Equations are
        resolved and factored as atoms with the binary atomic symbol \(=\);
        no other rule touches them.
\end{enumerate}
\emph{Output:} a derivation containing the empty clause certifies that \(S\)
is unsatisfiable.
\end{hbmethod}

\begin{hbtheorem}{thm:p02-w02b-p02-026-eq-resolution}{proved=annex:p02-w02b-p02-026-eq-resolution}{Soundness and refutational completeness of resolution with equality axioms}
Let \(S\) be a finite set of \(\Lang\)-clauses whose literals may contain
\(=\).  The following are equivalent.
\begin{enumerate}
  \item \(S\) is unsatisfiable under identity semantics.
  \item \(S^{\mathsf{eq}}\cup\mathrm{Eq}_S^{\mathsf{eq}}\) is unsatisfiable
        under the relational reading.
  \item Some finite set of ground instances of members of
        \(S\cup\mathrm{Eq}_S\), read as propositional clauses in which each
        distinct ground atom (ground equations included) is one
        propositional variable, is propositionally unsatisfiable.  Its ground
        terms can be chosen from the symbols of \(\Sigma_S\), together with
        one constant of \(\Lang\) when \(\Sigma_S\) contains none.
  \item Resolution with Factoring, as in
        Method~\ref{def:p02-w02b-p02-026-method}, derives the empty clause
        from \(S\cup\mathrm{Eq}_S\).
\end{enumerate}
The implication from (4) to (1) is the soundness of the method; the
implication from (1) to (4) is its refutational completeness.
\end{hbtheorem}

For an equality-free \(S\) the clauses (E1)--(E3) and the instances of (E4)
and (E5) can be dropped: Theorem~\ref{thm:fol-resolution} already applies.
The theorem above is the equality version of that result, obtained by
running the proved equality-free chain (Herbrand's theorem for clauses,
Theorem~\ref{thm:p02-w02-p02-028-clause-form}, and the lifting route of
Theorem~\ref{thm:p02-w02-p02-028-resolution}) on the translated set.  The
same quotient construction proves Herbrand's theorem with equality in its
validity form, Theorem~\ref{thm:p02-w02-p02-028-st-herbrand-eq}.

\begin{hbtheorem}{thm:p02-w02b-p02-026-bare}{proved=annex:p02-w02b-p02-026-bare}{What resolution without the axioms refutes}
Let \(S\) be a finite set of \(\Lang\)-clauses whose literals may contain
\(=\).  Resolution with Factoring derives the empty clause from \(S\) iff
\(S^{\mathsf{eq}}\) is unsatisfiable under the relational reading.  If
\(S^{\mathsf{eq}}\) is unsatisfiable, so is \(S\); the converse fails
(Corollary~\ref{thm:p02-w02b-p02-026-failures}).
\end{hbtheorem}

This is the exact content of the equality boundary of
\cref{sec:fol-cardinality-normalforms-04}: on clause sets with equality,
bare resolution is sound, and it is complete for the relational reading, but
not for identity.

\begin{hbexample}{def:p02-w02b-p02-026-example}{A refutation that needs a congruence axiom}
Let \(a,b\) be constants, \(f\) a unary function symbol and \(P\) a unary
relation symbol, and
\[
 S_0=\bigl\{\{a=b\},\ \{P(f(a))\},\ \{\neg P(f(b))\}\bigr\}.
\]
Under identity semantics \(S_0\) is unsatisfiable: \(a\) and \(b\) denote the
same element, so \(f(a)\) and \(f(b)\) do, and \(P(f(a))\) and
\(\neg P(f(b))\) cannot both hold.  Here \(\Sigma_{S_0}=\{a,b,f,P\}\), and the
clauses of \(\mathrm{Eq}_{S_0}\) used below are the function congruence for
\(f\) and the relation congruence for \(P\).  Number the clauses
\[
 \begin{array}{ll}
  1.\ \{a=b\} & 2.\ \{P(f(a))\}\\
  3.\ \{\neg P(f(b))\} & 4.\ \{\neg(x_1=y_1),\ f(x_1)=f(y_1)\}\\
  5.\ \{\neg(x_2=y_2),\ \neg P(x_2),\ P(y_2)\} &
 \end{array}
\]
(4 and 5 already renamed apart from each other and from the ground clauses).
\begin{enumerate}
  \item[6.] Resolve 1 and 4 on \(a=b\) and \(\neg(x_1=y_1)\) with most
        general unifier \(\{x_1\mapsto a,\ y_1\mapsto b\}\):
        \(\{f(a)=f(b)\}\).
  \item[7.] Resolve 2 and 5 on \(P(f(a))\) and \(\neg P(x_2)\) with
        \(\{x_2\mapsto f(a)\}\): \(\{\neg(f(a)=y_2),\ P(y_2)\}\).
  \item[8.] Resolve 7 and 3 on \(P(y_2)\) and \(\neg P(f(b))\) with
        \(\{y_2\mapsto f(b)\}\): \(\{\neg(f(a)=f(b))\}\).
  \item[9.] Resolve 6 and 8 on \(f(a)=f(b)\) and \(\neg(f(a)=f(b))\) with the
        identity substitution: the empty clause.
\end{enumerate}
By the implication from (4) to (1) of
Theorem~\ref{thm:p02-w02b-p02-026-eq-resolution} this derivation proves
\(S_0\) unsatisfiable.  Without clauses 4 and 5 no refutation exists
(Corollary~\ref{thm:p02-w02b-p02-026-failures}(1)): the atoms \(P(f(a))\)
and \(P(f(b))\) do not unify, and \(a=b\) has no complementary literal.
\end{hbexample}

\begin{hbcorollary}{thm:p02-w02b-p02-026-failures}{proved=annex:p02-w02b-p02-026-failures}{Two incomplete treatments of equality}
Let \(\Lang\) contain constants \(a,b\), a unary function symbol \(f\) and a
unary relation symbol \(P\).
\begin{enumerate}
  \item The set \(S_0\) of Example~\ref{def:p02-w02b-p02-026-example} is
        unsatisfiable, but no resolution derivation from \(S_0\) contains
        the empty clause.
  \item Let \(S_1=\bigl\{\{a=b\},\{P(a)\},\{\neg P(b)\}\bigr\}\) and
        \(\Sigma'=\{a,b\}\), which omits \(P\).  Then \(S_1\) is
        unsatisfiable, but no resolution derivation from
        \(S_1\cup\mathrm{Eq}(\Sigma')\) contains the empty clause.
\end{enumerate}
\end{hbcorollary}

Item~2 shows that the hypothesis on \(\Sigma\) in
Lemma~\ref{thm:p02-w02b-p02-026-quotient} cannot be dropped: a congruence
clause is needed for every relation symbol that occurs.

\begin{hbwarning}{def:p02-w02b-p02-026-limits}{Limits of the equality-axiom route}
\begin{enumerate}
  \item Theorem~\ref{thm:p02-w02b-p02-026-eq-resolution} asserts that a
        refutation exists when \(S\) is unsatisfiable.  It gives no bound on
        its length and says nothing about the behaviour of a search on a
        satisfiable \(S\).
  \item Nieuwenhuis and Rubio observe that resolution and factoring with
        these axioms tend to generate very many new clauses, most of them
        unnecessary, and present this as the reason for building equality
        into dedicated rules such as paramodulation
        \cite[\S1.1, p.~4]{NieuwenhuisRubio2001Typescript}.  The observation
        concerns search behaviour; it does not affect the soundness or
        completeness proved above.
  \item The clause set must be fixed before \(\mathrm{Eq}_S\) is formed.
        A symbol introduced later, for instance by Skolemization, without
        its congruence clause can make the method incomplete, as
        Corollary~\ref{thm:p02-w02b-p02-026-failures}(2) shows for a
        relation symbol.
\end{enumerate}
\end{hbwarning}

\subsection{The connection to superposition}
\label{sec:p02-w02b-p02-026-superposition}

The ordered-superposition calculus of the registered rule sheet
(Definition~\ref{def:p12-ordered-superposition}) removes the equality axioms
and builds identity into three rules: Superposition, which replaces an
instance of the larger side \(l\) of an equation \(l\approx r\) by \(r\) inside
another literal; Equality Resolution, which deletes a disequation
\(s\not\approx t\) whose sides unify; and Equality Factoring.  On that sheet
\(s\approx t\) is an equation under identity semantics, that is, the atom
\(s=t\) of this section, and \(s\not\approx t\) is \(\neg(s=t)\); the symbol
\(\approx\) is not the relational-reading symbol \(\mathsf{eq}\).  Every
literal is an equation or a disequation, so a relation atom
\(R(t_1,\ldots,t_n)\) has to be encoded as an equation before the rules apply.
For a term \(s\) and a position \(p\) of \(s\), \(s|_p\) is the subterm at
\(p\) and \(s[u]_p\) the result of replacing it by \(u\); the sheet writes
\(s[l']\) for a term with a distinguished subterm occurrence \(l'\) and
\(s[r]\) for the replacement at the same position.

\begin{hblemma}{thm:p02-w02b-p02-026-superposition-sound}{proved=annex:p02-w02b-p02-026-superposition-sound}{Soundness of the superposition rules}
Read the literals of the ordered-superposition rule sheet as above.  For each of the rules Superposition, Equality Resolution and
Equality Factoring of that sheet, every \(\Lang\)-structure that satisfies the
premises satisfies the conclusion.  The proof uses only that the stated
substitution is a unifier of the stated terms; the ordering, eligibility and
selection conditions are not used.  Consequently, if these rules derive the
empty clause from \(S\), then \(S\) is unsatisfiable.
\end{hblemma}

The converse, refutational completeness, depends on the ordering,
selection, redundancy and fair-saturation hypotheses of the sheet.  In this
handbook it is the cited
Theorem~\ref{thm:p12-superposition-completeness} and is not proved; the
equality-complete connection proved here is the one of
Theorem~\ref{thm:p02-w02b-p02-026-eq-resolution}.  The two routes are related
as follows.  Each Superposition step into a literal \(R(\ldots s\ldots)\)
does in one inference what resolution does with a congruence clause and a
chain of resolution steps, as in steps 6--8 of
Example~\ref{def:p02-w02b-p02-026-example}; the ordering restrictions then
forbid most of the replacements that the axioms would allow.  Nieuwenhuis and
Rubio trace this development from paramodulation, introduced by Robinson and
Wos, and Brand's result that the functional reflexivity axioms are not needed
\cite[\S1.1, p.~4]{NieuwenhuisRubio2001Typescript}, to superposition, the
restriction of paramodulation to inferences on the larger terms
\cite[\S1.2, p.~5]{NieuwenhuisRubio2001Typescript}, and to the completeness
of strict superposition for full first-order clauses with equality proved by
Bachmair and Ganzinger, whose proof the chapter presents in its Section~3
\cite[\S1.2, p.~6]{NieuwenhuisRubio2001Typescript}.

\emph{Sources.}  The clausal equality axioms and the satisfiability
equivalence are as in \cite[\S1.1, pp.~3--4]{NieuwenhuisRubio2001Typescript},
in the typescript version of the chapter; the resolution rules and their soundness
theorem are those of \cref{sec:fol-cardinality-normalforms-04}; the
equality-free completeness chain is
Theorems~\ref{thm:p02-w02-p02-028-clause-form}
and~\ref{thm:p02-w02-p02-028-resolution}.  The proofs of the results of this
section are written for this handbook.
